\documentclass[11pt]{article}
\usepackage[margin=1in]{geometry}
\usepackage{amsmath,amssymb,amsthm}
\usepackage{hyperref}
\usepackage{enumitem}

\newtheorem{theorem}{Theorem}
\newtheorem{lemma}{Lemma}
\newtheorem{corollary}{Corollary}

\newcommand{\Z}{\mathbb Z}
\newcommand{\ellinf}{\ell_\infty}
\newcommand{\R}{\mathbb R}
\newcommand{\diam}{\operatorname{diam}}

\title{A Sharp Cube Obstruction for the Cyclic-Metric Lower Bound in Creutz's Pu-Inequality Paper}
\author{Agent lane 3, \texttt{fa\_banach\_001}}
\date{2026-07-18}

\begin{document}
\maketitle

\section*{Source and Target Signal}

The source paper is Paul Creutz, \emph{Rigidity of the Pu inequality and quadratic isoperimetric constants of normed spaces}, arXiv:2004.01076, Rev. Mat. Iberoam. 38 (2022), 705--729.

Creutz proves that if $n\geq 2$ and $\mathcal A\geq\mathcal A^{ht}$ is an area functional, then
\[
  \operatorname{QIS}^{\mathcal A}(\mathbf{Ban}_n)
    = \left[\frac{1}{4\pi}, r_n^{\mathcal A}\right],
  \qquad
  r_n^{\mathcal A}<\frac{1}{2\pi},
\]
with $r_n^{\mathcal A}$ nondecreasing and converging to $1/(2\pi)$. The paper states that explicit values of $r_n^{\mathcal A}$ beyond the case $n=2$ remain open, and later states that for $n\geq 3$ the extremal spaces are completely open.

In Remark 4.9(1), the paper improves the basic cyclic-subset lower bound by observing that $S_{2n}$ embeds isometrically into $\R_\infty^n$, yielding
\[
  r_n^{ht}\geq C^{ht}(\R_\infty^n)
    \geq \left(1-\frac{2}{n}\right)\frac{1}{2\pi}.
\]
Here $S_m$ denotes $m$ equidistant cyclically ordered points in the intrinsic circle metric.

\section*{Open-Problem Transcription}

From the source introduction:

\begin{quote}
Explicit values of the optimal constants $r_n^{\mathcal A}$ beyond the aforementioned case $n=2$ remain open.
\end{quote}

From Example 4.7:

\begin{quote}
For $n\geq 3$ the question which spaces are extremal remains completely open.
\end{quote}

\section*{Result Claimed Here}

The following theorem does not compute the unknown constants $r_n^{\mathcal A}$. It gives a sharp obstruction to one natural way of trying to improve Creutz's explicit lower bound, namely by replacing the embedded cyclic metric $S_{2n}\subset \R_\infty^n$ with a larger cyclic metric inside the same cube.

\section*{Proof Intuition}

An isometric embedding into $\ell_\infty^N$ is the same as $N$ real-valued 1-Lipschitz coordinates such that every relevant distance is witnessed by at least one coordinate. For a long cycle, it is enough to count diameter pairs. On an even cycle, if a single coordinate witnesses one antipodal pair, its increments around the cycle are forced to be $+1$ on one semicircle and $-1$ on the other; that pattern can witness no other antipodal pair. On an odd cycle, a witnessing coordinate must contain a run of exactly $q$ consecutive $+1$ or $-1$ increments, where $m=2q+1$; because the total increment around the cycle is zero, there can be at most one such positive run and at most one such negative run. Thus one coordinate witnesses at most two diameter pairs. The matching upper bound is supplied by distance-to-landmark coordinates on a half-cycle.

\begin{theorem}
\label{thm:linf-dimension-cycle}
Let $m\geq 2$, and let $S_m=\Z/m\Z$ with cyclic graph metric
\[
  d(i,j)=\min\{|i-j|,\,m-|i-j|\}.
\]
The least integer $N$ for which $S_m$ embeds isometrically into $\ell_\infty^N$ is
\[
  N=\left\lceil \frac{m}{2}\right\rceil .
\]
\end{theorem}

\begin{proof}
Write $q=\lfloor m/2\rfloor$ and $p=\lceil m/2\rceil$.

For the upper bound, choose the consecutive landmark set
\[
  A=\{0,1,\ldots,p-1\}\subset \Z/m\Z
\]
and define
\[
  F:S_m\longrightarrow \ell_\infty(A),
  \qquad
  F(x)=(d(x,a))_{a\in A}.
\]
Every coordinate is 1-Lipschitz, hence
\[
  \|F(x)-F(y)\|_\infty\leq d(x,y).
\]
We show that equality holds. Fix $x\neq y$ and put $r=d(x,y)\leq q$. A landmark $a$ is good for the pair $(x,y)$ if
\[
  |d(a,x)-d(a,y)|=r.
\]
For a fixed pair at cyclic distance $r$, the bad landmarks form either two open arcs of lengths $r-1$ and $r-1$ when $m$ is even, or two arcs of lengths $r-1$ and $r$ when $m$ is odd. This is checked by rotating the pair to $(0,r)$ and evaluating the two displayed distance formulae directly. In the even case, if $r=q$, the only good landmarks are the two endpoints, and a set of $q$ consecutive vertices contains exactly one endpoint of every antipodal pair. If $r<q$, the two bad arcs each have length at most $q-2$, so a consecutive $q$-set cannot be contained in their union without crossing a good gap. In the odd case, the largest bad arc has length $r\leq q$, while $A$ has length $q+1$; again a consecutive $q+1$-set cannot lie inside the bad set, and when $r=q$ the complement of $A$ has only $q$ consecutive vertices, so it cannot contain both good endpoints of a diameter pair. Hence some $a\in A$ is good, and $F$ is an isometric embedding. This proves $N\leq p$.

For the lower bound, suppose $f:S_m\to\R$ is one coordinate of a 1-Lipschitz map. Put
\[
  \Delta_i=f(i+1)-f(i),\qquad i\in\Z/m\Z.
\]
Then $|\Delta_i|\leq 1$ and $\sum_i\Delta_i=0$.

First let $m=2q$ be even. There are $q$ unordered antipodal pairs $\{i,i+q\}$, $0\leq i<q$. If $f$ witnesses the distance of one such pair, then
\[
  |f(i+q)-f(i)|=q.
\]
Since this difference is the sum of $q$ increments, each of absolute value at most $1$, all $q$ increments on one semicircle must be equal to the same sign $\sigma\in\{\pm1\}$. Because the total sum of all $2q$ increments is zero, all increments on the complementary semicircle are equal to $-\sigma$. This sign pattern determines the antipodal pair. Thus one coordinate can witness at most one antipodal pair. An isometric embedding into $\ell_\infty^N$ must witness all $q$ antipodal pairs, so $N\geq q=m/2$.

Now let $m=2q+1$ be odd. The diameter is $q$, and there are $m$ unordered diameter pairs, for instance the pairs $\{i,i+q\}$ with $i\in\Z/m\Z$. If $f$ witnesses such a pair in the positive orientation, then
\[
  f(i+q)-f(i)=q,
\]
so the $q$ consecutive increments
\[
  \Delta_i,\Delta_{i+1},\ldots,\Delta_{i+q-1}
\]
are all equal to $+1$. Similarly, a witnessed pair in the opposite orientation gives a run of $q$ consecutive increments equal to $-1$.

There cannot be two distinct runs of $q$ consecutive $+1$ increments. Indeed, two distinct such cyclic runs on a cycle of length $2q+1$ would force at least $q+1$ increments to be $+1$; the remaining $q$ increments have sum at least $-q$, so the total sum would be positive, contradicting $\sum_i\Delta_i=0$. The same argument applies to $-1$ increments. Hence a single real coordinate can witness at most two diameter pairs: at most one through a $+1$ run and at most one through a $-1$ run.

Since the odd cycle has $2q+1=m$ diameter pairs and each coordinate can witness at most two of them, an isometric embedding into $\ell_\infty^N$ must have
\[
  2N\geq 2q+1,
  \qquad\text{hence}\qquad
  N\geq q+1=\left\lceil\frac{m}{2}\right\rceil .
\]
Combining the upper and lower bounds proves the theorem.
\end{proof}

\begin{corollary}
\label{cor:creutz-cube-sharp}
For each $n\geq 1$, the largest cyclic metric $S_m$ that can embed isometrically into $\R_\infty^n$ has $m=2n$ points.
Consequently, the cyclic-metric input in Creutz's Remark 4.9(1) is dimension-sharp within $\R_\infty^n$: no embedding $S_m\hookrightarrow \R_\infty^n$ with $m>2n$ can improve the displayed lower bound by the same route.
\end{corollary}

\section*{Verification Notes}

The proof is purely finite-dimensional and uses only the elementary coordinate characterization of $\ell_\infty^N$ embeddings. The included script \texttt{code/cycle\_linf\_dimension\_check.py} verifies, for small $m$, both the landmark upper-bound construction and the counting lower bound over all integer-valued 1-Lipschitz coordinate functions. These computations are sanity checks, not dependencies of the proof.

The packet claim is intentionally narrow. It does not exclude sharper lower bounds for $r_n^{ht}$ from other normed spaces, from non-cyclic finite subsets, or from a more direct filling-area argument.

\section*{Novelty and Literature Check}

Local run indexes were searched for arXiv:2004.01076 and the phrases \emph{Pu inequality}, \emph{quadratic isoperimetric constants}, \emph{cyclic metric}, \emph{cycle embedding}, and \emph{ell\_infty}. No duplicate packet for this exact cube-cycle obstruction was found. Web searches on 2026-07-18 for phrases including \emph{Rigidity of the Pu inequality quadratic isoperimetric constants}, \emph{quadratic isoperimetric constants normed spaces r\_n}, \emph{cycle metric normed space isometric embedding finite dimensional}, and \emph{cyclic metric normed space isometric embedding} found the source paper and adjacent metric-embedding material, but no later computation of the constants $r_n^{\mathcal A}$ and no target-local statement of Corollary~\ref{cor:creutz-cube-sharp}. The theorem may be known in metric-embedding folklore; the packet records it as a self-contained sharp obstruction relevant to the source target, not as a claimed full solution of the source open problem.

\section*{References}

\begin{itemize}[leftmargin=2em]
\item P. Creutz, \emph{Rigidity of the Pu inequality and quadratic isoperimetric constants of normed spaces}, Rev. Mat. Iberoam. 38 (2022), no. 3, 705--729; arXiv:2004.01076.
\end{itemize}

\section*{Human Review Recommendation}

Review Theorem~\ref{thm:linf-dimension-cycle}, with special attention to the odd-cycle counting step. If accepted, the result should be indexed as a partial/sharp-obstruction packet for arXiv:2004.01076: it proves that the paper's cube-based cyclic lower-bound construction cannot be improved by embedding a larger cyclic metric into the same $\ell_\infty^n$ target.

\end{document}
